\section*{Overview}

Convex hull trick is the DP optimization pattern for recurrences that can be rewritten as
\[
dp[i] = cost(i) + \min_j (m_j x_i + b_j)
\]
or the corresponding maximum version. The crucial step is not the data structure. It is the algebra that turns previous
states into lines and current states into query points.

\section*{Problem-Driven Motivation}

A recurring contest situation is:

\begin{itemize}
  \item \(n \le 2 \cdot 10^5\),
  \item the naive DP is \(O(n^2)\),
  \item after expanding squares or separating variables, every candidate \(j\) contributes something linear in the
  current variable \(x_i\).
\end{itemize}

At that point the real problem is no longer ``try every previous \(j\).'' It is ``maintain the lower envelope of many
lines.''

\section*{Recognition Pattern}

The strongest signals are:

\begin{itemize}
  \item a transition with \(\min\) or \(\max\) over previous states,
  \item algebraic expansion produces \(m_j x_i + b_j\),
  \item query points \(x_i\) are monotone, and maybe slopes \(m_j\) are monotone too,
  \item the DP bottleneck is the transition loop over all \(j\).
\end{itemize}

If the line insertion order or query order is arbitrary, the same derivation may still help, but the correct structure
becomes Li Chao tree instead of the deque hull in this note.

\section*{Derivation}

Suppose the transition is
\[
dp[i] = x_i^2 + C + \min_{j < i}\bigl(dp[j] + x_j^2 - 2x_i x_j\bigr).
\]

For fixed \(j\), the part depending on \(i\) is
\[
(-2x_j)\cdot x_i + (dp[j] + x_j^2).
\]

So state \(j\) becomes the line:
\[
y = m_j x + b_j,\qquad
m_j = -2x_j,\quad b_j = dp[j] + x_j^2.
\]

Now the optimization is geometric:

\begin{quote}
At query \(x_i\), find the line with minimum value.
\end{quote}

If lines are added in monotone slope order and query points are monotone, the lower hull can be stored in a deque.

\includegraphics[width=\textwidth]{figures/lower-hull.svg}

\section*{Worked Problem}

\paragraph{Problem.}
The coordinates \(x_0 < x_1 < \cdots < x_{n-1}\) are sorted. Compute
\[
dp[0] = 0,\qquad
dp[i] = x_i^2 + C + \min_{0 \le j < i}\bigl(dp[j] + x_j^2 - 2x_i x_j\bigr).
\]

\paragraph{Why naive fails.}
Trying all \(j < i\) for every \(i\) is \(O(n^2)\).

\paragraph{Why monotone CHT fits.}
Because the \(x_i\) are sorted:

\begin{itemize}
  \item query points \(x_i\) are increasing,
  \item slopes \(m_j = -2x_j\) are decreasing.
\end{itemize}

Those two monotonicities are exactly what the deque hull needs.

\paragraph{Deque invariant.}
The lines stored in the deque appear in the same order as the intervals of \(x\)-values where they are optimal.

\paragraph{Insertion rule.}
When a new line arrives, if the middle line of the last three is never optimal for any \(x\), remove it. That is the
``pop from the back'' step.

\paragraph{Query rule.}
When queries come with nondecreasing \(x\), if the second line is already better than the first at the current \(x\),
the first line will never become optimal again. That is the ``pop from the front'' step.

\section*{Implementation Reasoning}

The implementation uses exact integer comparisons with \texttt{\_\_int128}, not floating intersection points. That
avoids precision bugs and makes equal-slope handling explicit.

Important invariants:

\begin{itemize}
  \item slopes are inserted in monotone order,
  \item query points are processed in monotone order,
  \item every line removed from the deque is permanently useless under those assumptions.
\end{itemize}

The code below keeps the generic hull operations and also includes a direct solver for the quadratic DP above.

\section*{Correctness Intuition}

Each surviving line owns a contiguous interval of query \(x\)-values where it is best. Monotone slope insertion keeps
those ownership intervals in order. Once a line loses its entire interval because of a newer line, it can never come
back. Likewise, once the front line loses to the next line at the current monotone query point, all later query points
are even worse for that front line.

\section*{Complexity Analysis}

For the monotone deque version:

\begin{itemize}
  \item add line: amortized \(O(1)\),
  \item query: amortized \(O(1)\),
  \item total for \(n\) insertions and \(n\) queries: \(O(n)\),
  \item memory: \(O(n)\).
\end{itemize}

\section*{Common Pitfalls}

\begin{itemize}
  \item Using the deque hull when slopes or query points are not monotone.
  \item Deriving the line formula incorrectly even though the data structure is correct.
  \item Forgetting to handle equal slopes.
  \item Flipping inequalities when switching between minimum and maximum.
  \item Using floating intersections when exact integer arithmetic was available.
\end{itemize}

\section*{Variants and Failure Modes}

\begin{itemize}
  \item Li Chao tree handles arbitrary insertion/query order.
  \item Offline hulls can work when only one of the monotonicities is missing but queries can be resorted.
  \item If the recurrence does not reduce to linear functions, this trick is not applicable.
\end{itemize}

\section*{Practice Problems}

\begin{itemize}
  \item Quadratic DP transitions on sorted coordinates.
  \item Monotone slope/query problems where each previous state is a line.
  \item Problems that can be solved either by monotone CHT or by Li Chao after losing monotonicity.
\end{itemize}

\section*{References}

\begin{itemize}
  \item \href{https://cp-algorithms.com/geometry/convex_hull_trick.html}{cp-algorithms: Convex Hull Trick and Li Chao Tree}.
  \item \href{https://usaco.guide/CPH.pdf}{Competitive Programmer's Handbook}.
  \item \href{https://codeforces.com/catalog}{Codeforces Catalog}.
\end{itemize}
